\chapter{数学实验与编程资源}

数学家也有“实验室”，只是器材便宜得多：一支笔、一张纸、一个表格软件，或者一台能写几行代码的电脑。本附录介绍三类零门槛的实验工具，每类配一两个可以立刻上手的活动。先用一句话立规矩：\textbf{实验负责观察，证明负责下结论}——工具能帮你看得更多、猜得更准，但看得再多也不能替代证明。

\section{电子表格：最被低估的数学仪器}

不用装任何专门软件，电脑或手机里的表格程序就能做严肃探索。单元格里写公式，往下拖动，规律就自己浮出来。

\textbf{活动一：复利的威力。} 第一列填年份 $0$ 到 $30$，第二列第一年填本金 $1000$，下面每格写“上一格乘以 $1.05$”（年利率 $5\%$）。拖到底你会发现：$30$ 年后变成约 $4322$，是本金的 $4$ 倍多；更妙的是，大约每 $14$ 年翻一番。这就是指数增长，它解释了为什么“利滚利”听起来温和、时间一长却惊人。一个小技巧：把利率单独放在一个单元格里，让公式引用它——改一次利率，整条曲线跟着变。这个“把可变的量单独放出来”的习惯，就是数学里“参数”的思想。

\textbf{活动二：用筛子捞质数。} 第一列填 $1$ 到 $200$，第二列判断每个数是不是质数（试除它前面的质数即可）。数一数每一百里有多少个质数：$1$ 到 $100$ 有 $25$ 个，$101$ 到 $200$ 只剩 $21$ 个。你会亲眼看到“质数越往后越稀疏”——这个观察指向数论里一个大定理的方向。

\section{GeoGebra：会动的几何}

GeoGebra 是免费的动态几何软件（网页版即可使用）。它最大的本事是“拖动”：图形画好之后，拽住一个顶点晃动，所有依赖它的量实时跟着变。

\textbf{活动一：拖不垮的内角和。} 画任意三角形，让软件量出三个内角并显示它们的和。然后使劲拖动顶点，把三角形拉成胖的、扁的、尖的——和始终停在 $180^\circ$。这不是证明，但会让你真心觉得“这件事值得证明”。

\textbf{活动二：亲手逼近 $\pi$。} 画一个圆，再画出它的内接正多边形，把边数从 $6$ 加到 $12$、$24$、$48$，观察多边形周长与直径之比一步步逼近 $3.14159\dots$。拖动边数滑块时特别留意：多边形与圆之间的“缝隙”越来越小，却永远不会真正消失——这正是“逼近”二字最诚实的图景。两千多年前数学家们“割圆”求 $\pi$ 的思路，你在屏幕上几分钟就能重演。

\begin{shiyan}
没有电脑也能实验。拿一根绳子围着一个圆形杯口绕一圈，再量出杯口的直径，用绳子对着直径比一比：圆周的绳子长度总是直径的 $3$ 倍多一点。换瓶盖、换碗口、换盘子，倍数几乎不变——你刚刚“观察”到了 $\pi$ 的存在。
\end{shiyan}

\section{Python：让计算机替你数一万次}

Python 是一门语法接近英语的编程语言，网上有免费的在线运行环境，不用安装任何东西。对学数学来说，它的价值是“不知疲倦”：你想试一百万个例子，它不嫌烦。

\textbf{活动一：验证猜想。} 写一个小循环，对 $n$ 从 $1$ 到 $1000$ 检查“$n^2+n+41$ 是不是质数”。你会发现前几十个居然全是质数——这个著名的“质数制造机”公式能连对 $40$ 次。但程序继续跑下去，$n=40$ 时结果是 $41\times41$，轰然倒塌。做完这个实验，你会对“例子再多也不是证明”有肌肉记忆般的理解。

\textbf{活动二：撒点估 $\pi$。} 让程序在边长为 $1$ 的正方形里随机撒十万个点，数有多少落在内切圆里（到圆心距离小于半径）。圆面积与正方形面积之比是 $\dfrac{\pi}{4}$，所以“圆内点数除以总点数再乘 $4$”就是 $\pi$ 的估计值。点撒得越多，估计越稳——这是概率论“频率逼近概率”的直观演示。

\textbf{活动三：逼近 $\sqrt{2}$。} 让程序从区间 $[1,2]$ 开始，每次取中点，看中点的平方比 $2$ 大还是小，据此把区间砍掉一半，把 $\sqrt{2}$ 留在剩下那一半里。重复 $30$ 次，区间长度不到十亿分之一——$\sqrt{2}$ 被夹在一个极小的缝隙里。这就是“二分法”，也是极限思想的算法版：目标永远抓不到手里，但可以关进任意小的笼子。

做计算实验的基本功只有两条：改参数，看变化。把活动一里的 $1000$ 改成 $100000$，把公式换成 $n^2-n+41$，猜猜结果会怎样，再跑一遍验证——猜错往往比猜对更有收获。

\begin{fanli}
“电脑验证了一百万个例子，所以命题成立”——这句话永远是错的。上面那个质数公式连对 $40$ 次，第 $41$ 次就失败了；数学史上还有连对几亿次才出错的例子。程序只能排除假命题（找到反例）或增强信心，\textbf{确认}一个关于无穷多对象的命题，必须靠证明。实验是望远镜，不是判决书。
\end{fanli}

\section{去哪里找这些工具}

电子表格用任何常见办公软件都行；GeoGebra 搜索名字即可找到官网，网页直接打开；Python 可以搜“在线 Python 运行”找到免安装的网页环境。入门不需要看大部头教程：把本附录的活动照着做一遍，卡住时搜一句“电子表格 怎么写公式”或“Python for 循环”，边做边学最快。遇到报错别慌，报错信息是计算机在告诉你它哪里没听懂，这本身就是最实用的逻辑训练。

\begin{pandeng}
继续深入的关键词：数值计算（让计算机解方程）、蒙特卡洛方法（用随机撒点算面积和概率）、计算机辅助证明（著名的四色定理证明就靠计算机核验了大量情形，数学家们为此争论了多年：这算不算证明？）。这些话题在大学的计算数学与科学计算课程里等着你。
\end{pandeng}

\section*{思考题}

\textbf{观察题}
\begin{sikao}
\item 在活动一里把年利率改成 $10\%$，本金翻一番大约需要几年？提示：不必精算，观察第二列大约每隔几格翻倍，再和 $5\%$ 时的 $14$ 年比较。
\end{sikao}

\textbf{证明题}
\begin{sikao}
\item 证明：活动三每次“砍半”之后，$\sqrt{2}$ 一定落在留下的那个小区间里。提示：看中点的平方与 $2$ 谁大，分两种情况讨论即可。
\end{sikao}

\textbf{挑战题}
\begin{sikao}
\item 设计一个实验，估计“往画满平行线的纸上随机扔一根针，针压到线的概率”（布丰投针问题）。提示：先想清楚针落地时哪些量是随机的——位置是一个，还有别的吗？
\end{sikao}
